\documentclass[10pt,a4paper]{amsart}
\usepackage[margin=19mm]{geometry}
\usepackage{amsmath,amssymb,mathtools}
\usepackage[scr=rsfs,cal=euler]{mathalfa}
\usepackage{xfrac}
\usepackage[unicode,hidelinks,bookmarks=false]{hyperref}
\newcommand{\ie}{i.e.\ }
\newcommand{\cf}{cf.\ }
\newcommand{\resp}{respectively\ }
\newcommand{\Subsection}{\subsection}
\providecommand{\nobreakdash}{\nobreak\hbox{-}}
\setlength{\emergencystretch}{2em}
\multlinegap0pt
\usepackage{bm}
\usepackage{bbm}
\usepackage{dsfont}
\usepackage{xspace}
\usepackage{cancel}
\usepackage{mathtools}
\usepackage{amsthm}
\makeatletter
\@ifundefined{lemma}{\newtheorem{lemma}{Lemma}}{}
\@ifundefined{theorem}{\newtheorem{theorem}[lemma]{Theorem}}{}
\makeatother
\newcommand{\R}{\mathbb{R}}
\title[Killing Mean Curvature Solitons from Riemannian Submersions]{Killing Mean Curvature Solitons from Riemannian Submersions}
\author{Diego Artacho, Marie-Am\'elie Lawn, Miguel Ortega}
\date{Source: arXiv:2501.12998v2 (CC BY 4.0)}
\begin{document}
\maketitle

\noindent\emph{Source and licence.} Killing Mean Curvature Solitons from Riemannian Submersions. Version arXiv:2501.12998v2 (CC BY 4.0), distributed under the Creative Commons Attribution 4.0 International licence, as recorded in the arXiv licence metadata for this exact version and shown on the abstract page https://arxiv.org/abs/2501.12998v2. Full rights notice: licenses/arxiv-2501.12998-CC-BY-4.0.txt.\par\smallskip

\noindent\textit{Editorial scope.} Counted unit: the Theorem whose source label is \texttt{thm:pde} in the article (source0.tex lines 177--186, source label \texttt{thm:pde}), delivered together with the article's own complete original proof of it (source0.tex lines 188--289, one \texttt{proof} environment of 102 lines). No mathematics is written, repaired or re-derived: statement and proof are the article's own text line for line. Required context retained uncounted: eq:mcf3 (lines 115--117). Every definition, display and label of those lines is retained unchanged. Omitted from this package and not redistributed: 1--114, 118--176, 187--187, 290--636. Nothing inside a retained excerpt is omitted or altered: every retained body is delivered exactly as the article's lines are written, so no omission receipt of any kind accompanies this package. Citations: the retained excerpts carry no citation command at all, so no bibliography entry of the article is transcribed and no reference is redistributed. Reference adaptation: none was needed and none was made; every reference inside the delivered unit resolves inside it, so no unresolved reference or citation is produced. Printed slips kept literal: no printed word, symbol, hypothesis or number of the counted statement or of its proof was altered. Layout and class: the article's own class is \texttt{amsart} with packages including eurosym, amsfonts, amsmath, amssymb, amsthm, xcolor, latexsym, float, placeins, bm, inputenc, babel, a4wide, pdfsync; the wrapper is the lane's pinned \texttt{amsart} editorial wrapper at 10pt on A4, which restates the theorem-like environments the retained text needs and adds the article's own macro definitions that the retained text needs (\texttt{\textbackslash R}), with extra wrapper packages (bm, bbm, dsfont, xspace, cancel, mathtools). The delivered numbering is the unit's own. Assets: no retained span contains a figure environment, a graphics-inclusion command, a PDF-page inclusion, a table or a TikZ picture; no image file is copied into this package, the package declares no asset dependency and no image file is redistributed with it. Licence: version arXiv:2501.12998v2 of this article is distributed under the Creative Commons Attribution 4.0 International licence, as recorded in the arXiv licence metadata for this exact version and shown on the abstract page https://arxiv.org/abs/2501.12998v2. A full rights notice is delivered at licenses/arxiv-2501.12998-CC-BY-4.0.txt. Characters: every retained excerpt is pure ASCII, so the delivered engine, XeLaTeX with its default Latin Modern text font, reports no missing character.
\begin{equation}\label{eq:mcf3}
    H = g \left( K , \nu \right) \, .
\end{equation}

\begin{theorem}\label{thm:pde}
Let $(M,g_M)$ be an oriented Riemannian manifold, $\varphi \colon M \to (0,+\infty)$ a smooth function, $I \subseteq \R$ an interval and $(N,g) = (M \times I , g_M + \varphi dr^2)$, where $r$ is the usual coordinate in $\R$. Let $\mathcal{G}$ be the one-parameter family of isometries of $N$ given by $\mathcal{G}_t (x,r) = (x , r+t)$. Then, if $\Omega \subseteq M$ is open, the graph of a smooth function $u \colon \Omega \to I$ defines a $\mathcal{G}$-soliton on $N$ if and only if $u$ satisfies
\begin{equation}\label{eq:pde}
    \mathrm{div}\left(\frac{\nabla u}{W}\right)=\frac{1}{W}-\frac{1}{2W\varphi}g(\nabla \varphi,\nabla u) \, ,
\end{equation}
where the function $W \colon M \to (0,+\infty)$ is defined by
\begin{equation}\label{eq:defW}
    W = \sqrt{g_M(\nabla u , \nabla u) + \frac{1}{\varphi}} \, .
\end{equation}
\end{theorem}

\begin{proof}
The Killing vector field associated to the one-parameter family of isometries $\mathcal{G}$ is just the coordinate vector field $\partial_r$ corresponding to the interval $I$. It is easily verified that a unit normal vector field to $F$ is given by
\begin{equation*}
\nu = \frac{1}{W} \left( - \nabla u + \frac{1}{\varphi} \partial_r \right) \, .
\end{equation*}
Hence, the right-hand side of equation~\eqref{eq:mcf3} is
\begin{equation}\label{eq:Knu}
g \left( \partial_r , \nu \right) = \frac{1}{W} \, .
\end{equation}
Next, we need to compute the mean curvature $H$ of $F$. To this end, let $(e_1 , \dots , e_n)$ be a local $g_M$-orthonormal frame of $M$. The mean curvature $H$ is given by
\begin{equation*}
H = g \left( \vec{H} , \nu \right) = \sum_{i,j=1}^{n} \gamma^{ij} g \left( \nabla_{F_* e_i} F_* e_j , \nu \right) \, ,
\end{equation*}
where $\gamma^{ij}$ are the coefficients of the inverse of the matrix $\gamma_{ij} = g \left( F_* e_i , F_* e_j \right)$. By definition of $F$ as the graph of $u$, $F_* e_i = e_i + u_i \partial_r$, where $u_i = e_i(u)$. Note that $\nabla u = \sum_{i=1}^{n} u_i e_i$. Hence,
\begin{equation*}
    \gamma_{ij} = \delta_{ij} + \varphi u_i u_j \quad \text{and thus} \quad \gamma^{ij} = \delta_{ij} - \frac{1}{W^2} u_i u_j \, .
\end{equation*}
Moreover,
\begin{align*}
    \nabla_{F_* e_i} F_* e_j = \nabla_{e_i + u_i \partial_r} \left( e_j + u_j \partial_r \right) = \nabla_{e_i} e_j + u_{ij} \partial_r + u_j \nabla_{e_i} \partial_r + u_i \nabla_{\partial_r} e_j + u_i u_j \nabla_{\partial_r} \partial_r \, ,
\end{align*}
where $u_{ij} = e_i(u_j)$. Furthermore, using Koszul's formula and the obvious facts that $[\partial_r,e_i] = 0$ and $g(\partial_r,[e_i,e_j]) = 0$, we see that
\[
2g\left(\nabla_{\partial_r} e_i,e_j\right)
=\partial_r\left(g\left(e_i,e_j\right)\right)
+e_i(g\left(\partial_r,e_j\right))-e_j(g\left(\partial_r,e_i\right))
 =\partial_r\left( g\left(e_i,e_j\right)\right)=0 \, .
\]
and therefore $g(\nabla_{\partial_r} X,Y)=0$ for any vector fields $X,Y$ tangent to $M$ and independent of $r$. Hence, using also that $0 = \partial_r(\varphi) = \partial_r(g(\partial_r,\partial_r))=2g(\nabla_{\partial_r} \partial_r,\partial_r)$ and that the leaves of the foliation are totally geodesic, we obtain
\begin{align*}
   g& \left( \nabla_{F_* e_i} F_* e_j , \nu \right)  \\ =& g\left(\nabla_{e_i}e_j,\nu\right)
    +u_{ij}g\left(\partial_r,\nu\right)
    +u_j g\left(\nabla_{e_i} \partial_r,\nu\right)
    +u_i g\left(\nabla_{\partial_r} e_j,\nu\right)
    +u_i u_j g\left(\nabla_{\partial_r} \partial_r,\nu\right) \\
  =  & \frac{1}{\varphi W} g\left(\nabla_{e_i}e_j,\partial_r\right)
    -\frac{1}{W}g\left(\nabla_{e_i}e_j,\nabla u\right)  +\frac{u_{ij}}{W} \\
     & +\frac{u_j}{\varphi W}g\left(\nabla_{e_i}\partial_r,\partial_r\right)
    -\frac{u_j}{W}g\left(\nabla_{e_i}\partial_r,\nabla u\right)
     +\frac{u_i}{\varphi W}g\left(\nabla_{\partial_r}e_j,\partial_r\right) \\
    &-\frac{u_i}{W}g\left(\nabla_{\partial_r}e_j,\nabla u\right)
    + \frac{u_iu_j}{\varphi W}g(\nabla_{\partial_r}\partial_r,\partial_r) - \frac{u_iu_j}{W} g(\nabla_{\partial_r}\partial_r,\nabla u)\\
    = &\frac{-1}{W}g\left(\nabla_{e_i}e_j,\nabla u\right)+\frac{u_{ij}}{W}
    +\frac{u_j}{\varphi W}g\left(\nabla_{e_i}\partial_r,\partial_r\right)
    %-\frac{u_j}{W}A^L\left(e_i,\nabla u)\right)
    +\frac{u_j}{W}g\left(\partial_r,\nabla_{e_i}\nabla u\right)\\
    &+\frac{u_i}{\varphi W} g\left(\nabla_{\partial_r}e_j,\partial_r\right)-\frac{u_i}{W}g\left(\nabla_{\partial_r}e_j,\nabla u\right)-\frac{u_iu_j}{W}g\left(\nabla_{\partial_r} \partial_r,\nabla u\right)\\
    =& \frac{-1}{W}g\left(\nabla_{e_i}e_j,\nabla u\right)
    +\frac{u_{ij}}{W}+\frac{u_j}{\varphi W}g\left(\nabla_{e_i}\partial_r,\partial_r\right)
    +\frac{u_i}{\varphi W}g\left(\nabla_{\partial_r}e_j,\partial_r\right)
    -\frac{u_iu_j}{W}g\left(\nabla_{\partial_r} \partial_r,\nabla u\right).
    \end{align*}
As $\partial_r$ is a Killing vector field, we have that
\begin{equation}\label{eq:ci}
\varphi_i:=e_i(g(\partial_r,\partial_r))=2g\left(\nabla_{e_i}\partial_r,\partial_r\right)=-2g\left(\nabla_{\partial_r} \partial_r,e_i\right)=2g\left(\nabla_{\partial_r}e_i,\partial_r\right) \, .
\end{equation}
Moreover,
\[g\left(\nabla_{e_i}e_j,\nabla u\right)
=e_i(g(e_j,\nabla u))-g\left(\nabla_{e_i}\nabla u,e_j\right)=u_{ij}-g(\nabla_{e_i}\nabla u,e_j) \, , \]
obtaining
\begin{equation}
    g \left( \nabla_{F_* e_i} F_* e_j , \nu \right)=\frac{1}{W}g\left(\nabla_{e_i}\nabla u,e_j\right)+\frac{u_j\varphi_i}{2\varphi W}+\frac{u_i\varphi_j}{2\varphi W}-\frac{u_iu_j}{W}g\left(\nabla_{\partial_r} \partial_r,\nabla u\right).
\end{equation}
Putting everything together, we obtain
\begin{eqnarray*}
H&=& \sum_{i,j=1}^{n} \gamma^{ij} g \left( \nabla_{F_* e_i} F_* e_j , \nu \right)\\
&=&\frac{\mathrm{div}(\nabla u)}{W} -\sum^{n}_{i,j=1}\frac{1}{W^3}u_iu_jg\left(\nabla_{e_i}\nabla u,e_j\right)\\
&&+\frac{1}{W}\sum^{n}_{i,j=1}\delta_{ij}\left[-u_iu_jg\left(\nabla_{\partial_r} \partial_r,\nabla u\right)+\frac{u_j\varphi_i}{2\varphi}+\frac{u_i\varphi_j}{2\varphi}\right]\\
&&-\frac{1}{W^3}\sum^{n}_{i,j=1}\left[-u_i^2u_j^2g\left(\nabla_{\partial_r} \partial_r,\nabla u\right)+\frac{u_iu_j^2\varphi_i}{2\varphi}+\frac{u_i^2u_j\varphi_j}{2\varphi}\right]\\
&=&\frac{\mathrm{div}(\nabla u)}{W}
-\sum^{n}_{i,j=1}\frac{1}{W^3}u_iu_jg\left(\nabla_{e_i}\nabla u,e_j\right) +\frac{1}{W}\left[-\sum^{n}_{i=1}u_i^2g\left(\nabla_{\partial_r} \partial_r,\nabla u\right)+\sum^{n}_{i=1}\frac{u_i\varphi_i}{\varphi}\right]\\
&&-\frac{1}{W^3}\left[-|\nabla u|^4
g\left(\nabla_{\partial_r} \partial_r,\nabla u\right)+\frac{1}{\varphi}\sum^{n}_{i,j=1}u_i\varphi_iu_j^2\right]\\
&=&\frac{\mathrm{div}(\nabla u)}{W} -\sum^{n}_{i,j=1}\frac{1}{W^3}u_iu_jg\left(\nabla_{e_i}\nabla u,e_j\right)\\
&&+\frac{1}{W^3}\left[-W^2|\nabla u|^2g\left(\nabla_{\partial_r} \partial_r,\nabla u\right)+W^2\frac{g(\nabla \varphi,\nabla u)}{\varphi}\right]\\
&&-\frac{1}{W^3}\left[-|\nabla u|^4g\left(\nabla_{\partial_r} \partial_r,\nabla u\right)+\frac{g(\nabla \varphi,\nabla u)|\nabla u|^2}{\varphi}\right].
\end{eqnarray*}
Bearing in mind the definition of $W$~\eqref{eq:defW}, we get
\[ H = \frac{\mathrm{div}(\nabla u)}{W}-\sum^{n}_{i,j=1}\frac{1}{W^3}u_iu_jg(\nabla_{e_i}\nabla u,e_j)+\frac{1}{\varphi W^3}\left[-|\nabla u|^2g\left(\nabla_{\partial_r} \partial_r,\nabla u\right)+\frac{1}{\varphi}g(\nabla \varphi,\nabla u)\right].
\]
By~\eqref{eq:ci}, we see that $\nabla_{\partial_r} \partial_r=-\frac{1}{2}\nabla \varphi$. Hence,
\begin{equation} \label{eq:nH}
H=\frac{\mathrm{div}(\nabla u)}{W} -\sum^{n}_{i,j=1}\frac{1}{W^3}u_iu_jg(\nabla_{e_i}\nabla u,e_j)+\frac{1}{\varphi W^3}g(\nabla \varphi,\nabla u)\left(\frac{1}{2}|\nabla u|^2+\frac{1}{\varphi}\right).
\end{equation}
Notice now that, if $W_i = e_i (W)$,
\begin{eqnarray*}
\frac{\mathrm{div}(\nabla u)}{W}&=&
\mathrm{div}\left(\frac{\nabla u}{W}\right) -g\left(\nabla\left(\frac{1}{W}\right),\nabla u\right) =\mathrm{div}\left(\frac{\nabla u}{W}\right)+\sum_{i=1}^n\frac{W_i}{W^2}u_i\\
&=&\mathrm{div}\left(\frac{\nabla u}{W}\right)
+\sum_{i=1}^n  \frac{1}{2W^3}\left( -\frac{\varphi_i}{\varphi^2}+2g(\nabla_{e_i}\nabla u,\nabla u)\right)u_i\\
&=&\mathrm{div}\left(\frac{\nabla u}{W}\right) +\sum_{i,j=1}^n\frac{1}{W^3}u_iu_jg(\nabla_{e_i}\nabla u,e_j)-\frac{1}{2\varphi^2W^3}g(\nabla \varphi,\nabla u).
\end{eqnarray*}
By using equation~\eqref{eq:nH}, we obtain
\begin{eqnarray*}
H=\mathrm{div}\left(\frac{\nabla u}{W}\right)+\frac{1}{W^3}\left(\frac{1}{2\varphi}|\nabla u|^2+\frac{1}{2\varphi^2}\right)g(\nabla \varphi,\nabla u) \, ,
\end{eqnarray*}
and therefore
\begin{eqnarray}\label{H_Killinggraph}
H=\mathrm{div}\left(\frac{\nabla u}{W}\right)+\frac{1}{2W\varphi}g(\nabla \varphi,\nabla u) \, .
\end{eqnarray}
This, together with~\eqref{eq:Knu}, yields the result.
\end{proof}

\end{document}
